\section{The corrected six-contact height estimate}\label{sec:height}
Set
\[
 \Cset=\{0,-15,-20,-24,-30,-60\},\qquad
 F(u)=u(u+15)(u+20)(u+24)(u+30)(u+60),
\]
and $G(u)=7u^2+300u+3600$. For a finite prime set $S$ containing all
primes at most 324 and $u\notin\Cset$, define
\[
 M_S(u)=\prod_{\substack{p\notin S:\ v_p(u-c)>0\text{ odd}\\
                          \text{for some }c\in\Cset}}p.
\]
Negative valuations of $u$ do not enter the product. Put
\begin{align*}
 F_0&=324!,\qquad D=F_0\,3^{256F_0^4},\qquad G_0=116064000000,\\
 E_0&=10(DG_0)^2,\qquad Q_R=R^{RE_0},\\
 B_{\rm cm}&=D^2\,3^{21000D^2}+1,\\
 C_{\rm cm}&=2B_{\rm cm}+4096(1028+5D+84000D^2),\\
 H_6(S;R)&=2\cdot10^{13}\max(1,C_{\rm cm},Q_R^{17}),\\
 &\hspace{1em}R\ge\max(D,\max S,10^{12}).
\end{align*}
When the second argument is omitted, we use
$H_6(S)=H_6(S;\max(D,\max S,10^{12}))$.
The proof follows the corrected geometric route in S09, with the
constants of S08. It does not use the older alternative modularity route.

\begin{theorem}[Six-contact estimate]\label{thm:height}
For every such $S,R$ and every $u\in\Q\setminus\Cset$,
\[
 h(u)\le H_6(S;R)M_S(u)^{24}.
\]
More precisely the external arithmetic-geometric inputs below give
$h(u)\le C(S)M_S(u)^{24}$, where
\[
 C(S)=10^{13}\max((8D)^{288D},6^{100})\Bigl(\prod_{p\in S}p\Bigr)^{24}.
\]
\end{theorem}
\begin{proof}
Take $B=(2,-3)$, the maximal order $\mathcal O=\Z[i,(1+j)/2]$, and
$\mu=2j+ij$. With its positive sign choice, the form
$\Trd(\mu^{-1}x\overline y)$ has matrix
\[
 \begin{pmatrix}0&0&-1&1\\0&0&1&-2\\1&-1&0&0\\-1&2&0&0\end{pmatrix}
\]
and determinant one. This is the principal-polarization datum whose
normalizer quotient is $\mathbf P^1_j$, with branch values
$0,\infty,-16/27$ and indices $4,6,2$ \cite{babagranath,linyang}.
The full $\mathcal O$-linear level-three cover $Y_3$ has total degree
at most $4|\mathcal O/3\mathcal O|=324$, retaining all determinant
components. Here the normalizer quotient contributes four. Central
ideles act trivially after multiplying by a rational scalar and choosing
its sign to give residue one modulo three.

For an integral model apply \cite[Theorem 2.1, Corollary 2.3]{buzzard}
first at $V_1(5)$. This level has full determinant image, is maximal at
the division primes, and uses $5\ge4$ prime to the discriminant six.
Add full five-marking and full $\mathcal O$-linear three-marking, then
forget five-marking by a free finite quotient. The three-marking sheaf
is nonempty because an $\mathcal O_3$-stable lattice in the rank-one
$B_3$-module is a principal fractional left ideal. Freeness follows
because a polarized automorphism has finite order, while a finite-order
element congruent to one modulo three on $T_3$ is the identity.
The universal abelian surface descends. The auxiliary level contributes
no factor to the final degree.

Every component has genus at least two, by its negative orbifold Euler
characteristic, so the automorphisms extend to its smooth proper model.
At $p>324$ the action is tame. A stabilizer is cyclic and its tangent
character is faithful; a nontrivial fixed-point equation has invertible
linear term. Unique lifting prevents branch collisions. The three
ordered branch sections fix the integral $j$-coordinate, so no further
denominator prime is discarded.

Pull back by $J(u)=-57600F(u)^2/G(u)^6$. If
\[
 H(u)=(u^2-60u-1800)(13u^2+480u+3600)(11u^2+600u+7200),
\]
direct calculation gives $G^6-97200F^2=H^2$. The pullback indices are
$2,6,2$. Tame normalization replaces $e$ by $e/\gcd(e,a)$ and leaves
$2,1,1$. The discriminants and resultants have prime support at most 13,
so this persists at $p>324$. The normalized cover of $\mathbf P^1_u$
has degree at most 324, not $12\cdot324$, and is branched only at the
six contacts. Locally it has equation $u-c=z^2$ after unramified
extension. Positive even contact valuation is unramified; negative
valuation specializes to the unbranched point at infinity. This is
exactly the support specified by $M_S(u)$.

The splitting field of a fiber has degree at most $F_0$. The lattice
of ordered homomorphisms among its conjugate abelian surfaces has rank
at most $16F_0^2$. A finite subgroup of $\operatorname{GL}_r(\Z)$
injects modulo three. Adjoining its field of definition gives a normal
field $L$ of degree at most $D$, ramified only in
$S\cup\operatorname{supp}M_S(u)$. It defines all endomorphisms and all
ordered conjugacy homomorphisms.

The geometric endomorphism algebra is either $B$ or $M_2(k)$ for an
imaginary quadratic field $k$, by taking the right centralizer of the
complex structure. In the first case Skolem--Noether corrects conjugacy
isomorphisms to $B$-linear isogenies over $L$. These are in particular
compatible $\Q(\sqrt2)$-linear isogenies. Thus the normal-field,
endomorphism and compatible-isogeny hypotheses, including condition $(*)$,
of \cite[Eq. (7.1), Section 7.1.4, Proposition 7.13]{vkk} hold.
With $P_S=\prod_{p\in S}p$ and $M=M_S(u)$, this gives
\[
 h_F(A)\le(8D)^{288D}(P_SM)^{24}
\]
in the non-CM case.

In the CM case retain the degree bound $d_L=[L:\Q]\le D$.
The maximal-order reduction of \cite[Section 7.3.3, Eq. (7.17)]{vkk}
constructs, over the same field, an isogenous product
$A'=\prod A_i^{e_i}$ with maximal CM endomorphism rings. All geometric
endomorphisms were already defined over $L$, so the reduction needs
no further field extension. Its intermediate CM-type estimate gives
$h_F(A_i)\le6^{96}(P_SM)^{10}$, since the CM fields are unramified
outside the discriminant support of $L$. As $\sum e_i\le2$, put
$H_0=2\cdot6^{96}(P_SM)^{10}$; then $h_F(A')\le H_0$.
The isogeny estimate \cite[Eq. (7.13)]{vkk}, in dimension two, is
\[
 h_F(A)\le h_F(A')+\tfrac12\log\kappa(A'),\qquad
 \kappa(A')=\bigl(14^{32}d_L
             \max(h_F(A'),\log d_L,1)\bigr)^{16}.
\]
Since $\log14<3$, $D\ge2$, and $H_0\ge1$, it follows that
\[
 h_F(A)\le9H_0+16D+768\le(8D)^{288D}(P_SM)^{24}.
\]
For the last inequality use $18\cdot6^{96}<16^{64}$ and
$16D+768\le(8D)^4$, then absorb their sum in $(8D)^{65}$.
This argument bounds the original stable height while retaining the
field degree. It does not use the degree-free final conclusion of
Proposition 7.14, nor the historical alternative CM bound in S08.
The prescribed integer $C_{\rm cm}$ remains unchanged. Its former
degree-free interpretation is unnecessary: the corrected estimate is
absorbed by $Q_R^{17}$. Later comparisons between those integers are
numerical comparisons, not uses of the discarded height claim.

The compatible theta normalizations in \cite{linyang,lny} give
\[
 j=3779136\frac{(\prod z_m)^4}{(\sum z_m^8)^5},\qquad
 h(j)\le40h([z_m])+22,
\]
with ten even theta nulls; the other six vanish. The comparison in
\cite[Corollary 1.3(2)]{pazuki} yields
$h([z_m])\le10000\cdot2^{24}\max(1,h_F(A))$.
The change to the Deligne metric adds $\log\pi$ in dimension two and
has been included in this constant. A finite extension for a theta
structure changes neither absolute height. The coefficient norms of
$-57600F^2,G^6$ are below $4000^6$. The rational-root bound gives
$h(u)<h(j)+74$, after removing any initial power of $X$ or dropped
leading term. Since $40\cdot10000\cdot2^{24}+96<10^{13}$, this proves
the bound with $C(S)$.

Finally $P_S\le R^R$, $8D\le R^2$, and $6^{100}\le R^{576D}$.
Thus
\[
 C(S)\le10^{13}R^{576D+24R}\le10^{13}R^{600R}
 \le10^{13}Q_R^{17}\le H_6(S;R)/2.
\]
This also proves the assertion for every permitted larger $R$.
\end{proof}

The final value argument uses $S_{\rm ht}=\{p:p\le324\}$. Multiplication
uses $S_{\rm mul}=S_{\rm ht}\cup\mathcal H\cup\operatorname{supp}(c_*)$,
where $\mathcal H$ is the fixed catalogue. A common numerical $R$ may
bound both sets. It does not identify them.

\section{The pole predicate and its converse input}\label{sec:pole}
Let $S_0=\{2,3,5,7,13\}$. With input $B$ and eleven witnesses
$b_0,\ldots,b_5,e,z,r_1,r_2,w$, define
\begin{align*}
 m&=(1-2b_0)(1+b_1)(1+2b_2)(1+4b_3)(1+6b_4)(1+12b_5),\\
 A&=mB,\qquad H=A(A-B^2),\\
 T_1&=A(A-e)(A+3e),\qquad T_2=A(Az^2-eH)(Az^2+3eH),\\
 Q_{\rm pole}&=B\Bigl(\sum_{i=0}^5(b_i^2-b_i)^2+(eHzw-1)^2\\
 &\hspace{4em}+(r_1^2-T_1)^2+(r_2^2-T_2)^2\Bigr).
\end{align*}
Write $\Phi(B)$ for rational solubility. For $B\ne0$, $m$ is a signed
divisor of 2730, and the equations give points $(A,r_1)\in E_e$ and
$(Az^2,zr_2)\in E_{eH}$ with $eHz\ne0$, where
$E_\ell:y^2=x(x-\ell)(x+3\ell)$. The polynomial has total degree 75.

\begin{lemma}[One-sided obstruction]\label{lem:pole}
If $B\ne0$ and $\Phi(B)$, then $v_p(B)$ is not negative and odd for
any $p\notin S_0$.
\end{lemma}
\begin{proof}
If $k=v_p(B)<0$ were odd, then $v_p(A)=k$, $v_p(H)=3k$.
Exactly one of $e,eH$ has even valuation. Square scaling makes its
parameter a unit, preserving the parity of the abscissa's valuation.
For a unit parameter the cubic has valuation $v_p(x)$ if $v_p(x)>0$
and $3v_p(x)$ if $v_p(x)<0$. Its being a square forces even valuation.
But both points have abscissa in the squareclass $A$, a contradiction.
\end{proof}

\begin{hypothesis}[Pole completeness, $\mathsf{PC}$]\label{hyp:PC}
For every rational $B$ with square reduced denominator,
$0<B<1/2$ and $v_p(B)\ge4$ for every $p\in S_0$ imply $\Phi(B)$.
\end{hypothesis}
\begin{hypothesis}[Restricted converse, $\mathsf{RC}$]\label{hyp:RC}
For every positive squarefree $\ell\equiv11\pmod{12}$,
\[
 \operatorname{corank}_{\Z_2}\Sel_{2^\infty}(E_\ell/\Q)=1
 \quad\Longrightarrow\quad\ord_{s=1}L(E_\ell,s)=1.
\]
\end{hypothesis}
Here $c_4=208\ell^2$, $\Delta=2304\ell^6$, and $j=35152/9$.
The nonintegral rational $j$ shows these curves are non-CM. They are
different from the CM generator curve. Non-CM alone does not prove
Hypothesis~\ref{hyp:RC}.

\subsection{The exact descent reduction}
The signed divisors of 2730 realize all real, $2$-adic and $3$-adic
squareclass signatures. After $-1,2,3$ fix sign and valuation parities,
the remaining unit vectors of $5,7,13$ are
$(0,1,1),(1,0,0),(0,1,0)$ and form a basis. Choose positive $m$ making
$A=mB$ square at two and a nonsquare unit squareclass at three.
Let $\beta,\theta$ be the positive squarefree representatives of $A,H$.
Then $\beta\equiv17\pmod{24}$, $\theta\equiv1\pmod{24}$, with disjoint
supports, $\beta>1$, and $\theta=1$ permitted. Each is a local square
at the support of the other. At a prime of $\beta$, the relevant odd
valuation of $B$ is positive, so $1-B/m$ is a square unit. At a prime
of $\theta$, $v(B)$ must be zero, since otherwise $v(H)$ is even;
there $B/m\equiv1$, so $A$ is a square.

The simultaneous construction in S11--S12 gives $n,n\theta$ positive
squarefree and congruent to eleven modulo twelve, both with two-Selmer
dimension three and with extra Kummer class $(\beta,1)$. Its matrix
conditions can be stated explicitly. For support primes $q_i$, set
$s_i=[-1/q_i]$, $\lambda_i=[3/q_i]$, and
\[
 (Kv)_i=\sum_{j\ne i}[q_j/q_i](v_j+v_i).
\]
One arranges $\ker K=\langle\mathbf1,\mathbf1_{\operatorname{supp}\beta}\rangle$
and $\lambda\notin\operatorname{im}K$ on both supports. Add $d-1$
prime slots of type $(s,\lambda)=(0,0)$ and one slot $q_0$ of type
$(1,1)$, where $d=\omega(\beta)$. The cross-block has columns
$e_i+e_d$. After the rank-one correction $K^\sharp=K+\lambda a^t$,
deleting $q_0$ leaves
\[
 \begin{pmatrix}D_1&C_1\\C_1^t&0\end{pmatrix},
 \qquad\mathbf1^tD_1\mathbf1=1.
\]
This is nonsingular: $C_1^tx=0$ makes $x$ constant, and the first
equation forces that constant and then $y$ to vanish. Adding the
support of $\theta$ preserves this block. Free $q_0$-symbols toggle
the new diagonal entries; the determinant's full-product coefficient
is one. Thus a single Boolean assignment works on both supports.
Reciprocity, CRT and Dirichlet realize it by distinct auxiliary primes.

In Kummer coordinates $(x+3\ell,x-\ell)$, a positive-second-coordinate
Selmer representative is $(3^\delta\prod q_i^{u_i},\prod q_i^{v_i})$.
The necessary local equations are
\[
 Ku+\operatorname{diag}(\lambda)v=\delta\lambda,\quad
 (K+\operatorname{diag}(s))v=0,\quad(\lambda+s)^tv=0.
\]
The final equation is the condition at three. The matrix properties
force $v=0$, $\delta=0$, $u\in\ker K$, so there are at most eight
classes in total. The extra class is actually local at every place,
including two, and independent of the two torsion classes, giving
dimension three. This sufficiency check for the particular class is
essential; arbitrary parity conditions at two would not suffice.

\begin{lemma}[Dimension three gives corank one]\label{lem:corank}
For an elliptic curve with full rational two-torsion,
$\dim_{\F_2}\Sel_2=3$ implies $\operatorname{corank}_{\Z_2}\Sel_{2^\infty}=1$.
\end{lemma}
\begin{proof}
The finite Kummer sequence gives $r+\dim\Sha[2]=1$.
Write $\Sha[2^\infty]=(\Q_2/\Z_2)^t\oplus F$ with $F$ finite.
The Cassels--Tate pairing on the quotient by the maximal divisible
subgroup is nondegenerate and alternating for an elliptic curve
\cite{poonenstoll}. Its finite primary cyclic factors are paired,
so $\dim F[2]$ is even. Thus $F=0$ and $r+t=1$. The restriction of
the pairing to $F[2]$ need not be nondegenerate and is not used.
\end{proof}

Assuming $\mathsf{RC}$ for both $n$ and $n\theta$, the forward rank-one
theorem \cite{gz,kolyvagin} gives rank one and finite $\Sha$, hence
$\Sha[2]=0$.
The Kummer class is then represented by actual rational points. Their
abscissas have squareclass $\beta$, by the cubic equation. Write
$H=\theta c^2$ and rescale the second point to $E_{nH}$. If the first
abscissa is $Aq^2$, divide both parameters by $q^2$ and both coordinates
by $q^2,q^3$. The abscissas become $A,Az^2$, giving precisely the
eleven witnesses, with $w=(eHz)^{-1}$. If $\theta=1$, one can reuse
a point on the same curve. This proves $\mathsf{RC}\Rightarrow\mathsf{PC}$
with the stated descent input, not $\mathsf{PC}$ unconditionally.

\section{Multiplication and the finite false-triple guard}\label{sec:mult}
The actual multiplication template first forces all three input codes
to be members of $2E(K_0)$. These membership calls are conjuncts inside
the height certificate; they are not assumed caller preconditions.
Thus even untyped rational inputs, if accepted, have unique indices.
The template has 35,134 rational witnesses. Its construction is S03
and its mathematical specification is S13.

For indices $n_1,n_2,n_3$, each test group chooses $H=[h]P_0$,
$U_i=[z_i]P_0$, and sets $V_0=H-[c_*]P_0$,
$V_i=U_i-[c_*n_i]P_0$, where $c_*=4M_0$ is fixed. Use
\begin{align*}
 L(z)={}&16575z-6630d^2z^5+11050d^4z^9-25500d^6z^{13}\\
        &+68250d^8z^{17}-198900d^{10}z^{21}+612612d^{12}z^{25}.
\end{align*}
This is 16575 times the degree-25 truncation of the normalized formal
logarithm. The twelve tests are four inverse abscissas at $H,U_i$,
their four conjugate counterparts at $V_0,V_i$, three differences of
the corresponding $L$-ratios, and the product of the first two unprimed
ratios minus the third. Identity inverse abscissas are zero; all
denominators are guarded.

For a slope $t\notin\Cset\cup\{1\}$, put
\[
 R_0=\eta F(t)/(t-1)^8,\quad\eta^2=1,\quad
 a=1-R_0q^2\ne0,\quad R=\ell_0^2a^2R_0.
\]
Denote the twelve unscaled test values in group $g$ by $D_{gj}$.
Their normalized values are
\[
 C_{gj}=\frac{D_{gj}}{R_0^{k_j}(t-1)^{k_j+1}},\qquad
 k_j=\begin{cases}1,&1\le j\le8,\\25,&9\le j\le12.\end{cases}
\]
For six independent groups set
$W_g=\prod_j(1+RC_{gj}^2)$ and
$\Theta=\Tr\sum_gW_g^{-1}$, require them nonzero, and impose
$\Phi(R/\Theta^2)$. The pole input is nonzero, so the prefactor $B$
in $Q_{\rm pole}$ cannot make these constraints vacuous.

At a positive odd contact prime outside $S_{\rm mul}$, each factor
valuation is the minimum of zero and an odd integer. No cancellation
occurs. The trace and Lemma~\ref{lem:pole} force one common group and
place where the first eight tests have valuation at least one and
the last four at least 25. The first tests put the points in the
formal group. At least one of $h,h-c_*$ is a unit at the prime, so its full-log
ratio is integral. The truncated logarithm is the full logarithm times
a unit congruent to one in valuation at least 28. Its ratio is therefore
integral. The difference test transfers integrality to the other truncated
ratio, then to the other full-log ratio. Only now do the relative errors
give absolute congruences modulo $p^{28}$. Multiplying the two ratio congruences
by their integer denominators gives
$n_i\equiv L(z(U_i))/L(z(H))\pmod{p^{25}}$.
The product test then implies $p^{25}\mid n_1n_2-n_3$.

To specify the height certificate, write the abscissa of each input as
$x_{i0}+x_{i1}s$, using $x_{i0}=x_{i1}=0$ at the identity. Choose
$a_{i0},a_{i1},b_i\in\Q$ with $b_i\ne0$ and
$b_i x_{ij}=a_{ij}$. Express each of these nine rational quantities
as a sum of three slopes. Impose the preceding slope predicate on
all 27 slopes, with private witnesses, and impose membership on the
three input codes. This conjunction is the height certificate.
The multiplication predicate conjoins it with the fixed scalar
graphs and the inverse of the finite guard below.

\begin{proposition}[Finite-box deduction]\label{prop:box}
Assuming these local test identities and the height estimate for
$S_{\rm mul}$ with constant $A_{\rm ht}$, every false triple accepted
by the unguarded certificate has $\max_i|n_i|<T=(48A_{\rm ht})^{13}$.
\end{proposition}
\begin{proof}
Put $f=n_1n_2-n_3\ne0$. For each of the 27 slopes,
$M_{S_{\rm mul}}(t)^{25}\mid f$, so $h(t)\le A_{\rm ht}|f|^{24/25}$.
The nine rational homogeneous quantities for the three abscissas
are sums of three slopes. The sum and quotient height bounds give
$h(x_j)\le6A_{\rm ht}|f|^{24/25}+2\log3$ for each rational component.
The fixed coordinate change and the duplication bounds yield
\[
 2n_i^2\widehat h(P_0)\le12A_{\rm ht}|f|^{24/25}+B_0,
\]
where $2\widehat h(P_0)>1$ and
$B_0=4\log3+\tfrac32\log2+2\log(1+\sqrt2)+\tfrac13\log676<10$.
For $N=\max|n_i|$, $|f|\le2N^2$, so
$N^2<24A_{\rm ht}N^{48/25}+10$.
If $N\ge5$, this implies $N^{2/25}<48A_{\rm ht}$, hence
$N<(48A_{\rm ht})^{25/2}\le T$. The case $N\le4$ is immediate.
\end{proof}

For each false triple $(a,b,c)\in[-T,T]^3$, write its fifteen code
coordinates as $A_j/B_j$ with $B_j>0$, put $D_{abc}=\prod B_j$, and
form
\[
 g_{abc}(X)=\sum_{j=1}^{15}
       (D_{abc}X_j-(D_{abc}/B_j)A_j)^2.
\]
The guard requires the product of all these factors to have an inverse.
A false accepted triple makes its own factor zero. At a true triple
every factor is positive by injectivity of codes, and the inverse exists.
Thus the sound direction uses the generator and corrected all-$S$
height theorem but does not use pole completeness.

For completeness, S14 constructs, for every integer $A$ and bound $B_0$,
integers $u_1,u_2,v$ with $u_3=Av-u_1-u_2$ such that all 22 forms
$v,u_j-bv$, $b\in\Cset\cup\{1\}$, are signed catalogue multiples of
primes above $B_0$. At each such prime the conjugate reduction orders
have gcd four and are prime to that prime. The finite-complexity prime
asymptotic \cite{greentao,moebius,gtz} is used after fixing the catalogue
and initial CRT data, then the sieve cutoff $L$. Twenty-two catalogue
primes per unit class allow distinct choices for the twenty-two
occurrences. At one catalogue prime the distinguished zero plane avoids
at most $42p<p^2$ points; a lift modulo $p^2$ makes its divided value
different from zero and minus one. The fixed affine system is therefore
locally admissible. The asymptotic gives a lower count
$c_A(\log L)^{-44}X^3/(\log X)^{22}$ for fixed sieve cutoff $L$.
Each pair of forms has a two-by-two coefficient minor equal to
$\pm1$ or a nonzero contact difference of absolute value at most 61.
Thus the local conditions are nondegenerate uniformly in $A$;
other minors may involve $A$.
A Gaussian-pair sieve and the 21-form sieve on each nonzero fiber bound
the loss by $C_AX^3/(L(\log X)^{22})+o_A(X^3/(\log X)^{22})$.
Choose $L$ first and then $X$ with $A,L$ fixed. A uniform prime-asymptotic
error as $L$ varies is neither asserted nor needed. The corrected
all-orders inverse theorem is among the external inputs.

Apply this construction to the nine integral homogeneous quantities
of a true triple. Each slope has at most six odd-contact primes.
At one such prime with reduction orders $d_1,d_2$, solve
\[
 h\equiv0\pmod{\operatorname{lcm}(M_0,d_1)},\qquad
 h\equiv c_*\pmod{d_2}.
\]
Compatibility follows from
\[
 \gcd(\operatorname{lcm}(M_0,d_1),d_2)
                  \mid M_0\gcd(d_1,d_2)=4M_0=c_*.
\]
Choose a solution larger than $c_*$. Each contact prime uses its own
private group within its own slope. No common shift across different
primes is imposed, and selected prime values need not be distinct.
The saved library allocates six private groups per slope and twenty-seven
private slope calls. Set $U_i=[n_i]H$ and assign one group to each prime. Fixed local
clearing at the old primes can be checked as follows.
Put $S_b=S_{\rm ht}\cup\mathcal H$, $B_p=4$ on $S_{\rm ht}$ and
$B_p=1$ on the remaining catalogue primes, $C_p=14B_p$, and
$D_p=51C_p$. The selected integral forms give
$|v_p(R_0)|,|v_p(t-1)|\le C_p$, while the small formal subgroup
in the definition of $M_0$ gives $v_p(D_{gj})\ge0$ at $S_b$.
Consequently $v_p(C_{gj})\ge-D_p$. The fixed normalization has
$v_p(\ell_0)=E_p=52C_p+v_p(12)+4$.

Weak approximation chooses $q$ small at $S_b$, integral at the
selected primes, and arbitrarily close to a real root of
$1-R_0q^2$, without equality; choose $\eta$ to make $R_0>0$.
Then $a$ is a unit at the specified finite primes. At $p\in S_b$,
\[
 v_p(RC_{gj}^2)\ge 2E_p-C_p-2D_p>v_p(12),\qquad
 v_p(\Theta-12)>v_p(12),
\]
so $v_p(R/\Theta^2)\ge103C_p+8\ge4$.
At a selected odd-contact prime, $v_p(R)=1$. Each of the twelve
summands of the trace of the six $W_g^{-1}$ has residue zero or one;
the assigned group and place contributes one. Since $p>324$, their
sum is a unit and the pole input has valuation one. At every other
prime outside $S_b$, its valuation is even. Thus its reduced
denominator is square. In both real embeddings each $W_g$ is positive;
as $a\to0$, $W_g\to1$, $\Theta\to12$, and $R/\Theta^2\to0$
through positive values. This verifies exactly $\mathsf{PC}$'s
hypotheses, including all nonzero guards. Negative indices, zero indices and
coincident points are handled by the group-law branches. Therefore
the full multiplication interface remains conditional on
$\mathsf{PC}$, hence on the required restricted converse in this route.
